\chapter{\textit{complete}-Unabhängigkeit}
\label{ch:complete}

In diesem Kapitel wird das Entscheidungsproblem \(Ind_{co}\) untersucht. 
Zuerst erfolgt eine komplexitätstheoretische Einordnung, dann wird die eine \(labelling\)-basierte 
Kodierung der \textit{complete}-Semantik in aussagenlogischer Form vorgestellt. 
Dieses wird anschließend verwendet, um eine \textit{QBF}-Kodierung für \(Ind_{co}\) zu formulieren.

\section{Komplexität}

\section{Aussagenlogische Kodierung der \textit{complete}-Semantik}
Wir erinnern die Definition der \textit{complete}-Semantik. Sei \(F = (A,R)\) ein AF und sei \(S \subseteq A\) eine \textit{admissible}-Extension von \(F\).
\(S\) ist eine \textit{complete}-Extension genau dann, wenn für jedes Argument \(a \in A\), das \(S\) verteidigt, \(a \in S\) gilt.

\begin{Definition}
    \label{def:vars}
    Sei \(F = (A,R)\) ein AF. Für jedes Argument \(a \in A\) seien \(i_a,o_a,u_a\) aussagenlogische Variablen mit der Bedeutung
    \begin{align*}
        & i_a \text{ ist wahr genau dann, wenn } \lambda(a)=\mathit{in},\\
        & o_a \text{ ist wahr genau dann, wenn } \lambda(a)=\mathit{out},\\
        & u_a \text{ ist wahr genau dann, wenn } \lambda(a)=\mathit{und}.
    \end{align*}

    Für die \(k\)-te Kopie der Labelling-Variablen schreiben wir
    \begin{align*}
        & I^k=\{i_a^k\mid a\in A\},\\
        & O^k=\{o_a^k\mid a\in A\},\\
        & U^k=\{u_a^k\mid a\in A\}.
    \end{align*}
\end{Definition}

\begin{Definition}
    \label{def:pl-one}
    Sei \(F=(A,R)\) ein AF und seien \(I,O,U\) wie gewohnt definiert.
    Für \(a\in A\) definieren wir
    \[
        \mathrm{one}_a(I,O,U)
        :=
        (i_a \lor o_a \lor u_a)
        \land
        \Bigl(
            (\neg i_a \lor \neg o_a)
            \land
            (\neg i_a \lor \neg u_a)
            \land
            (\neg u_a \lor \neg o_a)
        \Bigr).
    \]
\end{Definition}

Die Formel \(\mathrm{one}_a(I,O,U)\) stellt sicher, dass dem Argument
\(a\) genau eines der drei Labels \(\mathit{in}\), \(\mathit{out}\)
oder \(\mathit{und}\) zugewiesen wird. Die erste Disjunktion fordert,
dass mindestens eines der drei Labels gilt. Die übrigen Konjunktionen
schließen aus, dass zwei verschiedene Labels gleichzeitig gelten.



\begin{Definition}
    \label{def:co-pl}
     Sei \(F = (A,R)\) ein AF und seien \(I, O, U\) wie zuvor definiert. 
     Eine aussagenlogische Kodierung eines \textit{complete}-Labellings lautet
     \[
     \mathrm{co}_F(I, O, U)
        :=
        \bigwedge_{a\in A}
        (
            \mathrm{one}_a(I,O,U)
        )
        \land 
        \bigwedge_{a\in A}
        (
            i_a
            \leftrightarrow
            \bigwedge_{(b,a)\in R} o_b
        )
        \land
        \bigwedge_{a\in A}
        (
            o_a
            \leftrightarrow
            \bigvee_{(b,a)\in R} i_b
        ).
     \]
\end{Definition}

\begin{Lemma}
    \label{lem:co-pl-correct}
    Sei \(F=(A,R)\) ein AF und sei \(\nu\) eine Belegung der Variablen
    \(I\cup O\cup U\). Falls für jedes \(a\in A\) genau eine der
    Variablen \(i_a,o_a,u_a\) unter \(\nu\) wahr ist, induziert \(\nu\)
    ein Labelling \(\lambda_\nu:A\to
    \{\mathit{in},\mathit{out},\mathit{und}\}\) durch
    \[
        \lambda_\nu(a)=
        \begin{cases}
            \mathit{in}, & \text{falls } \nu(i_a)=\top,\\
            \mathit{out}, & \text{falls } \nu(o_a)=\top,\\
            \mathit{und}, & \text{falls } \nu(u_a)=\top.
        \end{cases}
    \]
    Dann gilt:
    \[
        \nu\models \mathrm{co}_F(I,O,U)
        \quad\Longleftrightarrow\quad
        \lambda_\nu
        \text{ ist ein \textit{complete}-Labelling von } F.
    \]
\end{Lemma}

\begin{proof}

\end{proof}

\section{QBF-Kodierung der \textit{complete}-Unabhängigkeit}
Wir erinnern das Entscheidungsproblem \(Ind_{co}\).

\begin{Definition}
    Sei \(F=(A,R)\) ein AF und seien \(X,Y,Z\subseteq A\) paarweise disjunkte Mengen.
    Für \(k \in \{1,2,3\}\) sei \(\mathcal{L}^k = I^k \cup O^k \cup U^k\) 
    \begin{multline*}
    \Phi^{\mathrm{co}}_F(X,Y \mid Z)
    :=
    \forall \mathcal{L}^1 \mathcal{L}^2
    \exists \mathcal{L}^3
    \left(
        \left(
            \mathrm{co}_F(\mathcal{L}^1)
            \land
            \mathrm{co}_F(\mathcal{L}^2)
            \land
            \mathrm{Eq}^{1,2}_Z
        \right)
        \right. \\
        \left.
        \rightarrow
        \left(
            \mathrm{co}_F(\mathcal{L}^3)
            \land
            \mathrm{Eq}^{1,3}_X
            \land
            \mathrm{Eq}^{2,3}_Y
            \land
            \mathrm{Eq}^{1,3}_Z
        \right)
    \right).
    \end{multline*}
\end{Definition}

\begin{Lemma}
    Es gilt:
    \[
        \Phi^{\mathrm{co}}_F(X,Y \mid Z)\text{ ist wahr}
        \quad\Longleftrightarrow\quad
        X \bot_{\mathrm{co}} Y \mid Z.
    \]
\end{Lemma}